\section{Abstract Types Across a Boundary}
\label{sec:ch13-abstract-types-across-a-boundary}
\index{abstract types!across subgraphs|(}%

\index{abstract types!interface implementations in different subgraphs}%
The easy case is one you have already seen working, and it is worth naming
before the hard one. \code{Node} is declared in five of Mosaic's seven
subgraphs, and the four types implementing it are owned by four different
services: \code{Product} in Catalog, \code{Customer} in Accounts, \code{Review}
in Reviews, \code{Order} in Ordering. That composes, has composed since
chapter~\ref{ch:composition}, and the merged interface in
section~\ref{sec:ch13-giving-the-graph-a-node-field} is what came out.

So an interface whose implementations live in different subgraphs is not a
problem. Each subgraph declares the interface and its own implementations,
nobody declares anybody else's, and the composer merges the implementation
lists. That is the arrangement a team reaches for first, and it works.

Mosaic's other two abstract types do not cross anything: \code{SubmitReviewError}
and the \code{Error} interface belong to Reviews and nothing else declares them.
So a union whose members are owned by different services is a case this graph
has no instance of, and rather than inventing one in the storefront I have left
it as something this chapter did not exercise; the research note records it with
chapter~\ref{ch:inside-the-router-and-the-composer} as its owner.

\subsection*{The one that needs a directive}

\index{interfaceObject@\code{@interfaceObject}|(}%
What does not work by declaration alone is contributing a field to \emph{every}
implementation of an interface without knowing what the implementations are.
Chapter~\ref{ch:the-federation-model} described the directive for that from
Apollo's documentation. This is the smallest arrangement in which
\code{@interfaceObject} can be watched, and it lives under \code{samples/} in
the companion repository.

One subgraph owns a media library:

\begin{graphqlsdl}[fontsize=\footnotesize]
interface Media @key(fields: "id") { id: ID! title: String! }
type Book implements Media @key(fields: "id") { id: ID! title: String! pages: Int! }
type Film implements Media @key(fields: "id") { id: ID! title: String! runtimeMinutes: Int! }
\end{graphqlsdl}

\code{@key} on an \emph{interface} is the part to notice. It makes this an entity
interface rather than an ordinary abstract type, and it is what the other
subgraph attaches to. That other subgraph has never heard of books or films:

\begin{graphqlsdl}[fontsize=\footnotesize]
type Media @key(fields: "id") @interfaceObject {
  id: ID!
  averageRating: Float
  ratingCount: Int!
}
\end{graphqlsdl}

An interface on one side, an object on the other, and the directive is what tells
the composer they are the same type. Compose the pair and both new fields land on
the interface and on every implementation of it:

\begin{minted}[fontsize=\footnotesize,breaklines]{text}
Media  INTERFACE  id, title, averageRating, ratingCount
Book   OBJECT     id, title, pages, averageRating, ratingCount           implements Media
Film   OBJECT     id, title, runtimeMinutes, averageRating, ratingCount  implements Media
\end{minted}

Point a router at that and ask the library for a field the library has never
heard of:

\begin{minted}[fontsize=\scriptsize,breaklines]{text}
{ "__typename": "Book", "title": "The Mezzanine",  "averageRating": 4.666666666666667, "ratingCount": 3 }
{ "__typename": "Book", "title": "Pale Fire",      "averageRating": 5,                 "ratingCount": 2 }
{ "__typename": "Film", "title": "La Jetee",       "averageRating": 4.5,               "ratingCount": 4 }
{ "__typename": "Film", "title": "Sans Soleil",    "averageRating": null,              "ratingCount": 0 }
\end{minted}

\code{Sans Soleil} has no ratings row, and the field is nullable for exactly that
reason: a subgraph contributing to an interface cannot know which
implementations exist, so \enquote{I have nothing for this key} has to be an
ordinary answer rather than an error.

\subsection*{What the contributing subgraph is actually told}

\index{interfaceObject@\code{@interfaceObject}!the representation's typename}%
This is the mechanism, and it is stranger than the schema suggests.
Figure~\ref{fig:ch13-interface-object} has the shape of it.

\begin{figure}[htbp]
  \centering
  % What @interfaceObject does with a type name. Referenced from chapter 13. Bare
% tikzpicture; the figure environment, caption and label live at the call site.
%
% The point of the picture is the label on the middle arrow. The router reads
% keys off concrete types it knows about and sends them typed as the interface,
% so the subgraph on the right answers about a book without ever being told it
% is answering about a book.
\begin{tikzpicture}[
    font=\small,
    svc/.style={draw, rounded corners=2pt, minimum width=26mm,
                minimum height=15mm, align=center, font=\scriptsize},
    tp/.style={draw, rounded corners=2pt, minimum width=15mm,
               minimum height=5mm, align=center, font=\scriptsize},
    flow/.style={-{Stealth[length=2mm]}},
    note/.style={font=\scriptsize, gray, align=center, inner sep=1pt}
  ]

  % -- left: the owner --------------------------------------------------------
  \node[note, anchor=south] at (1.7,1.5) {\code{library}};
  \draw[rounded corners=2pt] (0.2,-1.9) rectangle (3.2,1.4);

  \node[tp] (media) at (1.7,0.9) {\code{interface Media}};
  \node[tp] (book)  at (1.7,-0.4) {\code{Book}};
  \node[tp] (film)  at (1.7,-1.3) {\code{Film}};

  \node[note, anchor=north, align=center] at (1.7,-2.1) {both entities,\\keyed on \code{id}};

  % -- right: the contributor -------------------------------------------------
  \node[note, anchor=south] at (8.6,1.5) {\code{ratings}};
  \draw[rounded corners=2pt] (7.1,-1.9) rectangle (10.1,1.4);

  \node[tp, minimum width=22mm] (io) at (8.6,0.5) {\code{type Media}\\\code{@interfaceObject}};
  \node[note, anchor=north, align=center] at (8.6,-0.8)
    {no \code{Book} here,\\no \code{Film},\\no way to list them};

  % -- the wire ---------------------------------------------------------------
  \draw[flow] (3.3,0.2) -- node[note, above, align=center]
    {keys read off\\\code{Book} and \code{Film}} node[note, below, align=center]
    {sent as\\\code{\_\_typename: Media}} (7.0,0.2);

\end{tikzpicture}

  \caption{The type name changes on the wire. The router reads keys off the
  concrete types the owning subgraph returned, and sends them to the
  contributing subgraph typed as the interface, which is the only type name that
  subgraph has. That is what lets it answer about a book without a \code{Book}
  type: it is never told that is what it is answering about.}
  \label{fig:ch13-interface-object}
\end{figure}

The sample's ratings subgraph writes down every reference-resolver call it
serves, in chapter~\ref{ch:entity-resolution-done-right}'s style. After the query
above:

\begin{minted}[fontsize=\footnotesize]{text}
_entities  Media/b1
_entities  Media/b2
_entities  Media/f1
_entities  Media/f2
\end{minted}

Four keys, all typed \code{Media}, and nothing about books or films. Now the
router's own plan for the same query:

\begin{minted}[fontsize=\scriptsize,breaklines]{text}
representations:
  @key  Book   fragment Key on Book  { __typename id }
  @key  Film   fragment Key on Film  { __typename id }
query: query($representations: [_Any!]!){ _entities(representations: $representations){
         ... on Media { averageRating } } }
\end{minted}

The router reads the keys as \code{Book} and \code{Film}, because that is what
the library returned, and sends them as \code{Media}, because that is the only
type the ratings subgraph has. The rewrite in the middle is the whole feature.
It is also why the promise holds for implementations that do not exist yet: add
a \code{Podcast} to the library tomorrow and it gets \code{averageRating} with no
change on the other side, because the other side was never told about the first
two either.

Narrowing the question narrows the work. Ask for the rating inside a fragment on
\code{Book} alone and the ratings log is two entries rather than four.

\subsection*{Three ways it goes wrong, and one of them is a crash}

\index{interfaceObject@\code{@interfaceObject}!forgetting the directive}%
Leave the directive off, so the contributing subgraph declares \code{Media} as an
ordinary object, and the composer is clear:

\begin{minted}[fontsize=\footnotesize,breaklines]{text}
"Media" is defined using incompatible types across subgraphs. It is defined as type "Interface" in subgraph "library"
but type "Object" in subgraph "ratings".
\end{minted}

That is the clearest message in this chapter and the only one that names both
subgraphs and the disagreement between them. Chapters~\ref{ch:composition}
and~\ref{ch:strangling-the-monolith} between them collected four different
mistakes that all report as a shareability error naming a type; this one reports
what actually happened.

\index{interfaceObject@\code{@interfaceObject}!an implementation with no key}%
Then the one that cost an afternoon. Give the interface a key, give one
implementation a key, and leave the other without one. There is no composition
error for this:

\begin{minted}[fontsize=\scriptsize,breaklines]{text}
Error: Fatal: Expected key "Film" to exist in the map "entityDataByTypeName".
    at invalidKeyFatalError (.../@wundergraph/composition/dist/errors/errors.js:433:12)
    at getOrThrowError (.../@wundergraph/composition/dist/utils/utils.js:28:49)
    at FederationFactory.handleEntityInterfaces (.../v1/federation/federation-factory.js:1337:68)
\end{minted}

wgc prints a box above that asking you to upgrade or open an issue, which is the
right thing for a tool to do when it has genuinely lost its footing, and it is
not much help when the cause is one missing directive in a schema you wrote a
minute ago. The rule to carry is short: \textbf{an entity interface needs every
implementation to be an entity, keyed the same way}. The composer does not
enforce it and cannot tell you when you break it.

This is committed as a case in \code{scripts/modeling-cases.mjs}, asserted by the
crash message, so a release that turns it into a proper error fails the gate
rather than quietly making this paragraph wrong.

\index{interfaceObject@\code{@interfaceObject}!an interface with no key}%
The third one I mention and did not chase. Give the contributing subgraph its key
and leave the owning \emph{interface} without one, and the pair composes. I never
pointed a router at that config, so I do not know what it serves, and the
research note records it as unmeasured with
chapter~\ref{ch:inside-the-router-and-the-composer} as its owner. What I will say
is that of the three, this is the one I would expect to reach production, because
it is the only one that produces no output at all.

\medskip
Interfaces and unions are structure, and structure is what a composer is good at
checking. The last two sections are about what happens as the structure runs out.

\index{interfaceObject@\code{@interfaceObject}|)}%
\index{abstract types!across subgraphs|)}%
